\documentclass[a4paper, 12pt]{article}

% Necessary Packages
\usepackage[utf8]{inputenc}
\usepackage[english]{babel}
\usepackage{amsmath}
\usepackage{amsfonts}
\usepackage{amssymb}
\usepackage{bm}       % For bold vectors (\bm or \mathbf)
\usepackage{geometry} % For better margins
\geometry{margin=2.5cm}
\usepackage{graphicx} % To include images if needed

% Header
\title{\textbf{IACV Homework Report}\\Geometric Rectification and 3D Reconstruction}
\author{Filippo Garofalo}


\begin{document}

\maketitle

\section{Introduction and Overview}
This report documents a complete pipeline for geometric rectification and metric 3D reconstruction from a single image of a church facade. The methodology follows a systematic approach grounded in projective geometry and computer vision theory:

\begin{enumerate}
    \item \textbf{Feature Extraction (Exercises 1-2):} We identify geometric primitives (lines and conics) in the image to determine vanishing points and the line at infinity.
    \item \textbf{Stratified Rectification (Exercise 3):} We progressively upgrade the geometric representation from projective to affine to metric space, restoring increasingly strong geometric properties.
    \item \textbf{Camera Calibration (Exercise 4):} Using orthogonal vanishing points, we determine the intrinsic parameters of the camera without explicit knowledge of the viewing geometry.
    \item \textbf{3D Reconstruction (Exercise 5):} We resolve scale ambiguity using a user-provided reference distance, enabling metric 3D reconstruction.
    \item \textbf{Structural Analysis (Exercise 6):} We determine the 3D position and orientation of the cylinder axis representing the vault structure.
\end{enumerate}

\section{Exercise 1}
\subsection*{Computation of the Line at Infinity}
The first step to recover the affine properties of the scene consists of identifying the image of the line at infinity ($\mathbf{l}_\infty$). In the 3D Euclidean world, parallel lines intersect at infinity. In perspective projection, these bundles of lines converge to a finite point on the image plane, known as the \textit{Vanishing Point} ($\mathbf{v}$).

By identifying two or more bundles of parallel lines in the scene (specifically, a bundle of vertical lines and a bundle of horizontal transversal lines), we compute the respective vanishing points $\mathbf{v}_{vert}$ and $\mathbf{v}_{trans}$ as the intersection of the fitted image lines.
The line at infinity $\mathbf{l}_\infty$ of the considered plane is the line passing through these two vanishing points:
\begin{equation}
    \mathbf{l}_\infty = \mathbf{v}_{vert} \times \mathbf{v}_{trans}
\end{equation}
This line encodes the geometric horizon of the plane and is fundamental for removing perspective distortion, enabling the transition from projective to affine space. In practice, the line is normalized by dividing by its third component ($\mathbf{l}_\infty := \mathbf{l}_\infty / \mathbf{l}_\infty[3]$) to enforce homogeneous coordinates and ensure numerical stability.
In the MATLAB code, 3 lines per direction were manually selected in order to retrieve the vanishing points for both the vertical and horizontal directions.

\section{Exercise 2}
\subsection*{Identification of the Cylinder Axis Vanishing Point}
The architecture of the church can be approximated as a cylinder whose main axis coincides with the depth direction of the nave.

To find the direction of this axis in the image space, we select lines parallel to the depth of the scene. The intersection of these lines generates the depth vanishing point, denoted as $\mathbf{v}_{depth}$ (or $\mathbf{v}_{axis}$). Ideally, this point represents the projection of the direction at infinity of the cylinder's axis onto the image plane.

In the context of a single-view reconstruction from a facade image, we have identified three mutually orthogonal vanishing points:
\begin{itemize}
    \item $\mathbf{v}_{vert}$: Vertical direction (up-down)
    \item $\mathbf{v}_{trans}$: Transversal horizontal direction (left-right)
    \item $\mathbf{v}_{depth}$ (or $\mathbf{v}_{axis}$): Depth direction (into the scene)
\end{itemize}
These three points form the basis for both the camera calibration in Exercise 4 and the 3D reconstruction in Exercises 5 and 6.

\section{Exercise 3}
\subsection*{Geometric Stratification: From Projective to Metric}

The goal of this exercise is to perform a full metric rectification of the scene. Instead of estimating a single transformation matrix directly, we follow a \textbf{stratified approach}. This involves upgrading the reconstruction level step-by-step:
\begin{center}
    Projective $\xrightarrow{\text{Step 1}}$ Affine $\xrightarrow{\text{Step 2}}$ Metric (Euclidean)
\end{center}
This order is necessary because each geometric stratum preserves specific invariants. We first need to restore parallelism (Affine) before we can reliably measure angles and shapes (Metric).

\subsubsection*{Step 1: Affine Rectification (Restoring Parallelism)}
The input image is a projective plane where parallel lines intersect at vanishing points. The first step, \textbf{Projective Transformation}, aims to map the line at infinity identified in Exercise 1 ($\mathbf{l}_\infty$) to its canonical position at infinity $\mathbf{l}_{\infty, canonical} = [0, 0, 1]^T$.

The basic projective rectification matrix $H_P$ is constructed as:
\begin{equation}
    H_P = \begin{bmatrix} 
    1 & 0 & 0 \\
    0 & 1 & 0 \\
    l_1 & l_2 & l_3 
    \end{bmatrix}
\end{equation}
where $\mathbf{l}_\infty = [l_1, l_2, l_3]^T$.

In the actual implementation, this matrix is scaled by its determinant to ensure proper numerical conditioning and improve the stability of the transformation:
\begin{equation}
    H_A = \det(H_P) \cdot H_P
\end{equation}

This scaled matrix $H_A$ is the one used for all subsequent transformations. Applying this transformation to the image eliminates projective distortion. In the resulting \textit{Affine Image}, parallel lines in the world appear parallel in the image. However, angles are still distorted, and circles appear as ellipses.

\subsubsection*{Step 2: Metric Rectification (Restoring Angles and Shapes)}
Once in the affine space, we can upgrade to a Metric (Euclidean) reconstruction. We observe that the diagonal arches of the vault (cross vault) are circular in the real world but appear as ellipses in the affine image.
To restore the circular shape (and consequently correct angles like orthogonality), we exploit the properties of conics.

The process involves the following substeps:
\begin{enumerate}
    \item \textbf{Conic Fitting:} We estimate the $3 \times 3$ symmetric matrix $C$ of the conic that best approximates the selected points of the diagonal arch in the original image.
    \item \textbf{Affine Transfer:} We transfer this conic into the affine space created in Step 1. Since the conic was defined in the original projective image space, we must apply the inverse transformation of $H_A$ to map it to the affine coordinate frame. Let $H_A^{-1}$ denote the inverse of the affine rectification matrix. The transformed conic $C_{aff}$ is given by:
    \begin{equation}
        C_{aff} = (H_A^{-1})^T \, C \, H_A^{-1}
    \end{equation}
    This ensures that the conic equation $\mathbf{x}^T C_{aff} \mathbf{x} = 0$ holds in affine coordinates.
    \item \textbf{Restoring Circularity:} In the affine space, the conic equation is $\mathbf{x}^T C_{aff} \mathbf{x} = 0$. The upper-left $2 \times 2$ submatrix of $C_{aff}$, denoted as $Q$, defines the shape (ellipticity) of the curve.
    
    To restore circularity, we apply two normalization steps:
    \begin{equation}
        Q \leftarrow \frac{Q}{Q(2,2)} \quad \text{(first normalization)}
    \end{equation}
    \begin{equation}
        Q \leftarrow \frac{Q + Q^T}{2} \quad \text{(symmetry enforcement)}
    \end{equation}
    
    For a circle in a Euclidean frame, this matrix should be proportional to the Identity matrix. Therefore, we seek a transformation $A$ such that the ellipse $Q$ is transformed into a circle. We decompose $Q$ via Cholesky factorization:
    \begin{equation}
        Q = A^T A
    \end{equation}
    The matrix $A$ represents the affine deformation. Its inverse represents the transformation needed to ``whiten'' the ellipse back into a circle.
\end{enumerate}

\subsubsection*{Final Transformation}
The metric update matrix $H_M$ is constructed using $A$:
\begin{equation}
    H_M = \begin{bmatrix} 
    A & \mathbf{0} \\
    \mathbf{0}^T & 1 
    \end{bmatrix}
\end{equation}
The total rectifying homography $H_{total}$ is the composition of the metric update applied after the affine rectification:
\begin{equation}
    H_{total} = H_M \, H_A
\end{equation}
Applying $H_{total}$ to the original image yields a result where both parallelism and metric properties (angles and length ratios) are correctly restored. Note that $H_A$ is the scaled affine rectification matrix (defined above), not the basic $H_P$.

\section{Exercise 4}
\subsection*{Camera Calibration (Analytical Method)}
To estimate the intrinsic parameters of the camera (Matrix $K$), we utilize the three vanishing points calculated previously: $\mathbf{v}_{vert}$, $\mathbf{v}_{trans}$, and $\mathbf{v}_{depth}$.

To simplify the problem and obtain a robust solution from a single view, we make two standard assumptions:
\begin{enumerate}
    \item \textbf{Zero Skew} ($s=0$): The pixels are rectangular.
    \item \textbf{Square Pixels} ($f_x = f_y = f$): The aspect ratio is 1.
\end{enumerate}

Under these assumptions, the \textit{Image of the Absolute Conic} (IAC), denoted as $\omega$, has a parameterized form with only 4 unknowns ($a, b, c, d$):
\begin{equation}
    \omega = \begin{bmatrix} a & 0 & b \\ 0 & a & c \\ b & c & d \end{bmatrix}
\end{equation}

\subsubsection*{Constraint Formulation}
For every pair of orthogonal vanishing points $(\mathbf{v}_i, \mathbf{v}_j)$, the orthogonality constraint $\mathbf{v}_i^T \omega \mathbf{v}_j = 0$ provides a linear equation in the unknowns $a, b, c, d$. 

Given three pairs of orthogonal VPs: (Vert-Depth), (Depth-Trans), and (Vert-Trans), we construct a linear system $\mathbf{A} \mathbf{x} = \mathbf{0}$ where $\mathbf{x} = [a, b, c, d]^T$ is the vector of unknowns. For each VP pair $(i,j)$ with coordinates $(x_i, y_i, w_i)$ and $(x_j, y_j, w_j)$, the constraint row is:
\begin{equation}
    (x_i x_j + y_i y_j), \quad (x_i w_j + w_i x_j), \quad (y_i w_j + w_i y_j), \quad (w_i w_j)
\end{equation}

This yields an over-determined system ($3 \times 4$) solved via Singular Value Decomposition: the solution is the last column of $V$ from the SVD decomposition $\mathbf{A} = \mathbf{U} \Sigma \mathbf{V}^T$.

\subsubsection*{Intrinsic Parameter Extraction}
Finally, the intrinsic parameters are extracted \textbf{algebraically} from $\omega$ using the relation $\omega = K^{-T} K^{-1}$. This yields:
\begin{itemize}
    \item Principal Point $x$-coordinate: $u_0 = -b/a$
    \item Principal Point $y$-coordinate: $v_0 = -c/a$
    \item Focal Length: $f = \sqrt{\frac{d}{a} - u_0^2 - v_0^2}$
\end{itemize}
The final Calibration Matrix $K$ is constructed as:
\begin{equation}
    K = \begin{bmatrix} f & 0 & u_0 \\ 0 & f & v_0 \\ 0 & 0 & 1 \end{bmatrix}
\end{equation}

\section{Exercise 5}
\subsection*{Reconstruction of 3D Coordinates}
Having obtained the calibration matrix $K$, we can back-project the 2D image points into 3D rays in the camera coordinate system. An image point $\mathbf{p} = [u, v, 1]^T$ defines a directional ray $\mathbf{d} = K^{-1}\mathbf{p}$, which is then normalized to unit length.

However, a single image suffers from scale ambiguity: the reconstructed points lie on correct rays but at unknown depths. To resolve this, we use a known reference measure: the user-specified unit distance between two base points.

The reconstruction procedure is as follows:
\begin{enumerate}
    \item We select two reference points on the image and record their 2D coordinates $(u_1, v_1)$ and $(u_2, v_2)$.
    \item We back-project them into normalized directional rays:
    \begin{equation}
        \mathbf{d}_{ref,i} = \frac{K^{-1}[u_i, v_i, 1]^T}{\|K^{-1}[u_i, v_i, 1]^T\|}
    \end{equation}
    \item We compute the distance between these two normalized rays: $\Delta d = \|\mathbf{d}_{ref,1} - \mathbf{d}_{ref,2}\|$
    \item We determine the scaling factor $\lambda$ such that the Euclidean distance equals 1 unit: $\lambda = 1.0 / \Delta d$
    \item We apply this global factor $\lambda$ to all back-projected rays of the 12 points on the diagonal arch, obtaining their metrically consistent 3D coordinates in the camera frame:
    \begin{equation}
        \mathbf{P}_i = \lambda \cdot \mathbf{d}_i
    \end{equation}
\end{enumerate}

\section{Exercise 6}
\subsection*{Localization of the Cylinder Axis}
The final objective is to determine the 3D position and orientation of the vault's axis. We define this as the \textbf{Cylinder Axis} passing through the apexes of the arch and aligned with the depth direction.

\subsubsection*{Position: Apex Localization}
The axis must pass through the highest point of the arch. The apex is identified as the \textbf{third point} from the initial 5-point ellipse selection (not necessarily the median, but the third element selected):
\begin{equation}
    \mathbf{p}_{apex,2D} = \text{pts\_ellisse}(\cdot, 3)
\end{equation}
This point is back-projected into 3D space:
\begin{equation}
    \mathbf{P}_{apex} = \lambda \cdot \frac{K^{-1} \mathbf{p}_{apex,2D}}{\|K^{-1} \mathbf{p}_{apex,2D}\|}
\end{equation}
where $\lambda$ is the scale factor computed in Exercise 5.

\subsubsection*{Direction: Level Constraint and Orthogonality}
The axis of the nave must be parallel to the ground plane (level). Relying solely on the raw depth vanishing point $\mathbf{v}_{depth}$ might introduce a slight tilt due to noise or measurement errors.

To ensure mathematical orthogonality, we construct a Rotation Matrix $R$ from the three orthogonal vanishing points. This is accomplished by:
\begin{enumerate}
    \item Back-projecting each vanishing point into a directional vector using $K^{-1}$
    \item Normalizing each vector to unit length
    \item Stacking them as columns: $R_{raw} = [\mathbf{d}_x \,|\, \mathbf{d}_y \,|\, \mathbf{d}_z]$
    \item Enforcing orthogonality via SVD: $R_{raw} = U \, \Sigma \, V^T \Rightarrow R = U V^T$ (this is the nearest orthogonal matrix in the Frobenius norm sense)
\end{enumerate}
The third column $\mathbf{r}_3$ of $R$ represents the depth direction that is guaranteed to be orthogonal to both the vertical ($\mathbf{r}_2$) and transversal ($\mathbf{r}_1$) directions. We ensure correct orientation by checking the sign with respect to the original vanishing point.

Thus, the cylinder axis is defined parametrically as:
\begin{equation}
    \text{Axis}(t) = \mathbf{P}_{apex} + t \cdot \mathbf{r}_3, \quad t \in \mathbb{R}
\end{equation}

\end{document}